% 37_body.tex — 산출물 ㊲ 제품-프로젝트 경계 합의서 · SLA/OLA/UC · 배상 상한 협상 메모 (빈 템플릿 / 완성 예시 공용 본문) — gen_product_templates.py 가 Product_specs.py 에서 생성
% 드라이버: 37_경계합의서SLA배상_빈템플릿.tex (\wbblanktrue) / 37_경계합의서SLA배상_완성예시.tex (\wbblankfalse — 워크북 §X.5 확정 후)
\documentclass[10pt,a4paper]{article}
\usepackage[a4paper,margin=16mm,top=14mm,bottom=20mm]{geometry}
\def\DocNum{37}\def\DocDay{8}\def\DocIdx{2}
\def\DocTitle{제품-프로젝트 경계 합의서 · SLA/OLA/UC · 배상 상한 협상 메모}
\def\DocSub{Product/Project Boundary, SLA Mapping \& Liability Negotiation}
\def\DocVariant{\B{빈 템플릿}{딥게이지 완성 예시}}
\usepackage{wbstyle}
\def\DocCirc{㊲}   % newunicodechar(㉛~㊵ 폴백) 뒤에 정의해야 활성 문자로 읽힌다
\begin{document}
\docheader
 {\Bg{[딥게이지 · 온담F\&B · 온담푸드시스템]}{딥게이지 · 온담F\&B홀딩스 · 온담푸드시스템}}
 {\Bg{[v1.x / 2028-05]}{v1.0 / 2028-05-24 (D+814)}}
 {\Bg{[작성자 — CPO·법무]}{윤하경 CPO · 법무 자문\rd{7mm}}}
 {\Bg{[서명자 — 양측]}{강민수 · 나영선 · 조성태\rd{7mm}}}

%% ------------------------------------------------------------------ 1
\secbar{1. 당사자와 문서 정보}
\guide{3자(딥게이지 / 온담F\&B홀딩스 / 온담푸드시스템). 상위 계약·SOW·SLA와의 관계.}
{\footnotesize
\begin{tabularx}{\linewidth}{|L{36mm}|Y|}\hline
\hd{항목} & \hd{내용} \\ \hline
\ifwbblank
\rd{8mm}당사자·대표자 &  \\ \hline
\rd{8mm}상위 계약 / 부속 문서 &  \\ \hline
\rd{8mm}버전 / 기준일 &  \\ \hline
\rd{8mm}적용 범위(검사 대상 4종 중) &  \\ \hline
\else
당사자·대표자 & ㈜딥게이지(강민수 CEO) / ㈜온담F\&B홀딩스(서명권자 나영선 — 정미란 CFO 결재 후 효력) / ㈜온담푸드시스템(조성태 대표) \\ \hline
상위 계약 / 부속 문서 & 마스터 계약(21억) 부속 A(경계) · B(SLA) · C(배상) · D(이관). 유상 SOW는 부속 A 참조. 충돌 시 마스터 → 이 문서 → SOW \\ \hline
버전 / 기준일 & v1.0 / 2028-05-24 · 3자 회의(강민수 · 나영선 · 조성태, 박세연·최영주 배석) \\ \hline
적용 범위(검사 대상 4종 중) & 4종 전부 — ① 이천 냉장(한동석) ② 61점 HACCP(나영선) ③ 입고 검수(나영선 + 구매) ④ 안성 납품 검사(조성태 법인) \\ \hline
\fi
\end{tabularx}}

%% ------------------------------------------------------------------ 2
\secbar{2. 경계 3영역 — 항목 배치}
\guide{제품 라인 / 유상 커스텀 / 고객 자체 개발. ㊱ 3분류를 옮겨 적고 경계 항목의 판정 기준을 명시.}
{\footnotesize
\begin{tabularx}{\linewidth}{|L{30mm}|Y|Y|}\hline
\hd{영역} & \hd{포함 항목(C-ID)} & \hd{판정 기준} \\ \hline
\ifwbblank
\rd{12mm}제품 라인 &  &  \\ \hline
\rd{12mm}유상 커스텀 &  &  \\ \hline
\rd{12mm}고객 자체 개발(온담 IT본부) &  &  \\ \hline
\rd{12mm}경계 미확정 — 협의 대상 &  &  \\ \hline
\else
제품 라인 & C-02·09·10·12·14·15·23·24·29·32·33·39·46·47 (14) & 자동차 58사에게도 가치 + 커널(판정 대행 없음·감시 아님) 무충돌 — 로드맵 편입·무상·시점 미약속 \\ \hline
유상 커스텀 & C-01·03·04·08·16·17·18·19·21·26·27·28·31·34·35·38·40·44·45 (19) & 온담 전용이되 모듈로 격리 가능 — 별도 SOW 59 mm × 1,200만 = 7.08억 \\ \hline
고객 자체 개발(온담 IT본부) & C-05·06·11·20·25·36·37·43 (8) & 마스터·법정 DB·IoT·P1 포털·DR — 온담 데이터 거버넌스·인프라 소관. 딥게이지는 인터페이스 사양 제공 \\ \hline
경계 미확정 — 협의 대상 & C-07(→ C-08 축소, 유상) · C-22(측정 후 재논의, 보류) · C-30·41(커널 충돌, 자체) · C-42(계약 별지) & 회의에서 확정 → 미확정 0 \\ \hline
\fi
\end{tabularx}}

%% ------------------------------------------------------------------ 3
\secbar{3. 소유권·재사용권·유지보수 책임}
\guide{영역별 — 코드·모델·데이터·문서의 소유권, 딥게이지의 재사용권(다른 고객에 판매), 유지보수·업그레이드 책임과 비용.}
{\footnotesize
\begin{tabularx}{\linewidth}{|L{30mm}|Y|Y|Y|}\hline
\hd{영역} & \hd{소유권} & \hd{재사용권} & \hd{유지보수 책임·비용} \\ \hline
\ifwbblank
\rd{11mm}제품 라인 &  &  &  \\ \hline
\rd{11mm}유상 커스텀 &  &  &  \\ \hline
\rd{11mm}고객 자체 개발 &  &  &  \\ \hline
\rd{11mm}학습 데이터 &  &  &  \\ \hline
\else
제품 라인 & 딥게이지(코드·모델·문서) & 해당 없음 — 제품 & 딥게이지 · 라이선스 포함 \\ \hline
유상 커스텀 & 코드·문서 온담 / 모델 가중치 딥게이지 & 딥게이지 — 온담 식별 정보를 제외한 재사용권(다른 고객에 판매 가능) & 딥게이지 · 연 SOW 금액의 18\% \\ \hline
고객 자체 개발 & 온담 & 해당 없음 & 온담 · 인터페이스 변경 60일 전 통지 \\ \hline
학습 데이터 & 원본 온담 / 가공·라벨 딥게이지 사용권 & 딥게이지 — 모델 학습 사용, 제3자 제공 불가 & 측정 세트 80건의 라벨 판정자 = 정본 판정자(항목 7) \\ \hline
\fi
\end{tabularx}}

%% ------------------------------------------------------------------ 4
\landscapeon
\secbar{4. SLA / OLA / UC 정합성 매핑}
\guide{SLA 99.5\% · 온프렘 · 릴리스 분기 1회 승인 후. 각 SLA 항목을 받치는 OLA(딥게이지 내부)·UC(GPU 벤더·클라우드)의 합이 SLA보다 강한가. PM 트랙 ⑰과 같은 표를 공급자 쪽에서.}
{\footnotesize
\begin{tabularx}{\linewidth}{|Y|L{22mm}|Y|Y|L{30mm}|L{16mm}|}\hline
\hd{SLA 항목} & \hd{목표} & \hd{OLA(내부)} & \hd{UC(외부)} & \hd{사슬 합 vs SLA} & \hd{성립?} \\ \hline
\ifwbblank
\rd{10mm} &  &  &  &  &  \\ \hline
\rd{10mm} &  &  &  &  &  \\ \hline
\rd{10mm} &  &  &  &  &  \\ \hline
\rd{10mm} &  &  &  &  &  \\ \hline
\rd{10mm} &  &  &  &  &  \\ \hline
\rd{10mm} &  &  &  &  &  \\ \hline
\else
가용성(월) & 99.5\% = 219.2분 & 운영 0.7 FTE 평일 09-18 · 대기 0.5 FTE & GPU 벤더 NBD 교체 · 온담 IT 1차 대응 & 2대 이중화 → 단일 무중단 · 이중 장애 NBD 24-72h > 219분 & 조건부(이중화 전제) \\ \hline
P1 복구 & 4시간 & 주간 2h · 야간 대기 4h & 온담 IT 야간 1차 30분 & 야간 0.5 + 4 = 4.5h > 4h & 불성립 → 최선 노력 또는 온담 OLA \\ \hline
릴리스 & 분기 1회, 온담 승인 후 & 스테이징 검증 2주 & 온담 IT 승인 3주(E-14) & 승인 지연 = 모델 갱신 불가 & 성립 — 승인 지연 기간 판정 오류 SLA 제외 \\ \hline
추론 지연(P95) & 5초 & 2-pass 4.3초(V-05) & GPU 2대 부하 분산 & 이중 장애 시 8.6초 & 조건부 \\ \hline
데이터 정정 & 2영업일 & CS 온담 담당 1명(1일) & 온담 정본 판정자 확인 1일 & 1 + 1 = 2 & 성립 — 판정자 지정 전제 \\ \hline
AI 판정 정확도 & FN ≤ 3.0\%(측정 세트 한정) & eval 분기 1회 재측정 & 온담 라벨 판정자 & 온담 4세트 미측정 — 세트 없으면 행 자체가 없다 & 임시 기준(㊵ 항목 4 전) · 크레딧, 배상 아님 \\ \hline
\fi
\end{tabularx}}
\landscapeoff

%% ------------------------------------------------------------------ 5
\secbar{5. 배상 상한 협상 메모}
\guide{딥게이지 10\% vs 온담 100\%(연 7억). AI 판정 오류를 SLA 대상에 넣을 것인가 — 넣는다면 어떤 오류(FN만? 법정 기준 판정?)를, 어떤 상한으로. 물리적 하한 계산(구축비·원가를 검산한 뒤). 4구간 제안.}
{\footnotesize
\begin{tabularx}{\linewidth}{|L{40mm}|Y|}\hline
\hd{항목} & \hd{내용} \\ \hline
\ifwbblank
\rd{11mm}온담 요구 / 우리 초안 &  \\ \hline
\rd{11mm}AI 판정 오류의 SLA 편입 여부와 범위 &  \\ \hline
\rd{11mm}법정 기준(④) 판정의 분리 조항 &  \\ \hline
\rd{11mm}배상 상한의 물리적 하한 계산 &  \\ \hline
\rd{11mm}4구간 제안(10 / 25 / 50 / 100\%)과 각 구간의 조건 &  \\ \hline
\rd{11mm}권고안 &  \\ \hline
\else
온담 요구 / 우리 초안 & 온담 100\%(연 청구 7.0억 전액) / 딥게이지 10\%(㉘ 항목 6 책임 한계의 연장 — 세영 162만 선례) \\ \hline
AI 판정 오류의 SLA 편입 여부와 범위 & 편입 — 단 측정 세트(80건)에서 측정된 FN에 한정, 위반 시 서비스 크레딧(무상 기간)이며 배상 대상 아님. 자동판정 기본 OFF, 최종 판정은 온담 검사원(커널 6-1) \\ \hline
법정 기준(④) 판정의 분리 조항 & 'GAUGE는 법정 기준 대비 판정 보조 정보를 제공하며 법정 기준 준수 판정과 그 결과의 책임은 온담에 있다. 법정 기준 DB의 탑재·갱신은 온담 소관(C-11).' \\ \hline
배상 상한의 물리적 하한 계산 & 구축 실소요 6.12(3.00 아님, F-14) · 운영 2.22 · 제품화 2.70 → 1년차 현금 −4.04억 · 3년 공헌 5.52억(SOW 별도) · 현금 7.52억. 100\% = 7.0억 = 공헌의 127\% = 현금의 93\%. 50\% = 3.5(63/47\%) · 25\% = 1.75(32/23\%) · 10\% = 0.7(13/9\%) \\ \hline
4구간 제안(10 / 25 / 50 / 100\%)과 각 구간의 조건 & 10\% 일반 SLA 위반(크레딧, 연 상한) / 25\% 이관 검증 6종 미통과·output 미달(㊵ 항목 6) / 50\% 보안 사고(유출·감사 로그 훼손) / 100\% 고의·중과실만. AI 판정 오류·법정 판정은 어느 구간에도 없다 \\ \hline
권고안 & 25\% 구간을 열어 주고(온담이 두려워하는 것은 이관 실패 — 11년치 41,000권) 100\%의 조건(고의·중과실)을 지킨다. 초안 10\% 고수는 관계를, 100\% 수용은 회사를 잃는다 — 상한은 지키고 서비스로 관계를 산다 \\ \hline
\fi
\end{tabularx}}

%% ------------------------------------------------------------------ 6
\secbar{6. 보안성 검토 218문항 대응 계획}
\guide{ISMS-P 미인증 · 5영업일 · 제218문항(사고 RCA·재발 방지 증빙 → ㉞ 항목 8). 답할 수 있는 것 / 조건부 / 못 하는 것.}
{\footnotesize
\begin{tabularx}{\linewidth}{|L{30mm}|Y|Y|L{40mm}|}\hline
\hd{구분} & \hd{문항 범주} & \hd{대응} & \hd{증빙} \\ \hline
\ifwbblank
\rd{11mm}즉시 답변 가능 &  &  &  \\ \hline
\rd{11mm}조건부(기한·비용) &  &  &  \\ \hline
\rd{11mm}답변 불가 — 대안 &  &  &  \\ \hline
\else
즉시 답변 가능 & 131 — 조직·정책·접근통제·감사 로그·전송 암호화·개인정보 처리·사고 이력(제218문항) & 정책 문서 12종 회신 · 제218문항은 ㉞ v1.1 전문 + 재발 방지 6건 증빙 + 감시 대시보드 리포트 6개월치 & ㉞ 항목 8 증빙 열 그대로 \\ \hline
조건부(기한·비용) & 62 — 침투 테스트·SBOM·키 관리·백업 복구 훈련·SSDLC & 외부 모의해킹 1,800만 · 4주(C-45) · SBOM(C-34) · 복구 훈련 기록 분기 & 계약 후 8주 내 회신 \\ \hline
답변 불가 — 대안 & 25 — ISMS-P·ISO 27001·SOC 2 인증서 · DR 센터 · 24/7 SOC & 2029-Q4 ISMS-P 획득 계획 · 온프렘이므로 인프라·DR은 온담 통제(C-37) · SOC는 온담 NMS 연동(C-38) & 계획서 + 온담 인프라 통제 확인서 \\ \hline
\fi
\end{tabularx}}

%% ------------------------------------------------------------------ 7
\secbar{7. 데이터 이관 조항}
\guide{정본 판정자(나영선/한동석/조성태 중 계약서에 지정) · 6종 검증 · 코드 매핑 커버리지 61.4\%의 처리(38.6\%는 온담의 마스터 정비 후) · 구축비 3억의 범위.}
{\footnotesize
\begin{tabularx}{\linewidth}{|L{40mm}|Y|L{44mm}|}\hline
\hd{항목} & \hd{문구} & \hd{리스크} \\ \hline
\ifwbblank
\rd{10mm}정본 판정자 &  &  \\ \hline
\rd{10mm}검증 6종과 통과 기준 &  &  \\ \hline
\rd{10mm}매핑 불가분(38.6\%)의 처리 &  &  \\ \hline
\rd{10mm}구축비 범위·추가 비용 조건 &  &  \\ \hline
\rd{10mm}일정·책임 &  &  \\ \hline
\else
정본 판정자 & '이관 데이터의 정본 판정자는 나영선 품질안전본부장으로 한다. 이천(한동석)·안성(조성태)은 위임 판정자, 상충 시 나영선 최종.' & 정본은 미지정(§5.5) — 지정 없이는 검증 6종의 판정이 없다 \\ \hline
검증 6종과 통과 기준 & 건수 100\% / 합계 오차 0 / 키 유일성 100\% / 참조 정합성 ≥ 99.5\% / 코드 매핑 커버리지 ≥ 61.4\%(1차) → 마스터 정비 후 95\% / 표본 대조 300건 오류 ≤ 0.5\% & 95\%를 1차 기준으로 쓰면 이관 실패 조항 \\ \hline
매핑 불가분(38.6\%)의 처리 & '매핑 불가 항목은 임시 식별자로 이관하고 판정 초안 대상에서 제외. 온담 마스터 정비(C-05) 완료 후 재매핑, 재매핑은 별도 견적. 마스터 정비 기한 계약 후 6개월(온담 의무).' & 기한 없으면 38.6\%가 영원히 임시 \\ \hline
구축비 범위·추가 비용 조건 & '구축비 3.00억은 커버리지 61.4\% 기준 이관(종이 11,200권 · 엑셀 4,100개)을 범위로 한다. 범위 밖(전체 41,000권·재매핑)은 별도.' & 3.00으로 6.12어치를 약속하면 1년차 −4.04가 −7.16 \\ \hline
일정·책임 & 서명 후 이관 6개월(PM 상주) · 검증 6종 통과 → 인수(정본 판정자 서명) · 미통과 시 25\% 구간 & 온담 CFO 결재 3주가 이관 시작을 미룬다 \\ \hline
\fi
\end{tabularx}}

%% ------------------------------------------------------------------ 8
\secbar{8. 합의 서명}
\guide{3자 서명. 서명 전 확인: 경계 미확정 항목이 0인가.}
{\footnotesize
\begin{tabularx}{\linewidth}{|Y|Y|Y|Y|}\hline
\hd{딥게이지 강민수} & \hd{온담F\&B 나영선} & \hd{온담푸드시스템 조성태} & \hd{(입회)} \\ \hline
\ifwbblank
\rd{16mm} & & &  \\ \hline
\else
\rd{16mm} 강민수 (서명) & 나영선 (서명 · CFO 결재 조건) & 조성태 (서명) & 박세연 · 최영주 (배석) — 미확정 0건 확인 \\ \hline
\fi
\rd{7mm}일자 & & &  \\ \hline
\end{tabularx}}

\end{document}